\section{Как это считают программы}
\label{sec:codes}

Все программы вычисляют одну величину --- производную \eqref{eq:xi-trace} или эквивалентное ей среднее \eqref{eq:xi-ham}, --- но в разных моделях элементов и разными путями: конечной разностью по $\delta$, продифференцированной картой с энергией-параметром или нормальной формой. Различия результатов определяются не методом, а тем, какие члены главы~\ref{sec:terms} и~\ref{sec:edge} содержит модель диполя. В этой главе разобраны BELL, COSY INFINITY, MAD-X/PTC и OptiM.

\subsection{Карты Тейлора и Ли (BELL)}

Одно-оборотное отображение строится как композиция карт элементов; карта диполя получается из гамильтониана тела \eqref{eq:H} (как ряд Тейлора или как Ли-фактор $\exp(:f:)$), без краевых отображений. Переменная --- канонический $\delta=\delta_p$. Хроматичность берётся прямо по определению: при фиксированном $\delta$ находится неподвижная точка \eqref{eq:fixed}, вычисляются дробные числа колебаний якобиана \eqref{eq:nu-trace}, и $\nu(\delta)$ дифференцируется центральной разностью. Высшие $\nu^{(j)}$ получаются так же из карты высокого порядка. Модель содержит члены A, C, E, S и не содержит F. Нормальная форма 6D (Драгт--Финн) рассчитана на кольцо с ВЧ и для $\xi$ в 4D не используется.

\subsection{Дифференциальная алгебра (COSY INFINITY)}

COSY строит карты в дифференциальной алгебре (DA): каждая координата --- усечённый многочлен от начальных координат и, возможно, от параметров~\cite{Berz1999,COSYManual}. Команда \texttt{OV}~$n$~$n_v$~$n_p$ задаёт порядок $n$, число фазовых пар $n_v$ и число параметров $n_p$. Переменные COSY --- $(x,a,y,b,l,\delta_K)$ с $a=p_x/p_0$, $b=p_y/p_0$~\cite{ValetovPhD}, поэтому сырые хроматичности относятся к $\delta_K$ и переводятся по \eqref{eq:xi-conv}.

\paragraph{Энергия как параметр: \texttt{OV 3 2 1}.}
Кинетическая энергия референсной частицы задаётся как $K_s\cdot\mathrm{PARA}(1)$, т.\,е. $K=K_s(1+\varepsilon)$ с DA-переменной $\varepsilon\equiv\delta_K$. Карта четырёхмерна, но её коэффициенты --- многочлены от $\varepsilon$. Нормальная форма (\texttt{TP}) поворачивает две поперечные пары, и числа колебаний получаются как функции параметра: $\nu(\varepsilon)=\nu_0+(\mathrm{d}\nu/\mathrm{d}\delta_K)\,\varepsilon+\dots$ Это в точности определение \eqref{eq:nu-series}: неподвижная точка при данном $\varepsilon$ находится в DA автоматически, и второе слагаемое \eqref{eq:dM} учтено~\cite{ValetovPhD}. Условие корректности --- поля магнитов должны оставаться фиксированными при изменении $\varepsilon$; если бы градиенты масштабировались с жёсткостью, хроматичность была бы тождественно нулевой.

\paragraph{Энергия как координата: \texttt{OV 3 3 0}.}
Карта шестимерна, $\delta_K$ --- шестая координата, и в ней есть коэффициенты $(x|\delta_K)$, $(x|x\delta_K)$, $(a|a\delta_K)$ и т.\,д. Хроматичность извлекается методом неподвижной точки~\cite{ValetovPhD}: в блоке $(x,a,\delta_K)$ решается $\lambda=(M-I_H)^{-1}e_3$, $I_H=\operatorname{diag}(1,1,0)$, что даёт $(D,D')$; карта пересчитывается около $\mathring z=(D,D',\dots,0,\delta)\delta$, и \eqref{eq:xi-trace} записывается явно через коэффициенты первого и второго порядков (приложение~\ref{app:map}). Результат тот же, что при \texttt{OV 3 2 1}: обе постановки дифференцируют одно и то же семейство отображений.

\paragraph{Почему \texttt{TP} не работает на 6D-карте без ВЧ.}
Нормальная форма приводит линейную часть к трём поворотам. Без ВЧ продольный блок
\begin{equation}
\label{eq:jordan}
M_{l\delta}=\begin{pmatrix}1&(l|\delta_K)\\0&1\end{pmatrix}
\end{equation}
имеет двукратное собственное значение $1$ и не диагонализуется --- это жорданова клетка. Поворота нет, продольного действия нет, и нормальная форма 6D не существует. С ВЧ блок становится поворотом на $2\pi\nu_s$, и \texttt{TP} работает, но $\delta$ уже не параметр, а фазовая переменная: линейная хроматичность в нормальной форме не появляется (глава~\ref{sec:rf}).

\paragraph{Краевые поля: \texttt{FR}.}
Режим \texttt{FR 0} --- жёсткий край без краевого поля: диполь описывается только телом, член F отсутствует, и результат совпадает с BELL. Режимы \texttt{FR 1}--\texttt{FR 3} включают краевое поле со спадом Энге, масштабированным апертурой; \texttt{FR 3} интегрирует уравнения движения в поле, удовлетворяющем уравнениям Максвелла, на всём протяжении края. Эти режимы содержат $B_s$ и, следовательно, член F; при апертуре $\to0$ они стремятся к пределу жёсткого края \eqref{eq:xiF}, а при конечной апертуре отличаются от него на поправки $O(g)$ раздела~\ref{sec:edge}. Для \texttt{FR 1}, \texttt{FR 2} нужна таблица \texttt{SYSCA.DAT}~\cite{COSYManual}.

\subsection{MAD-X и PTC}

\paragraph{MAD-X \texttt{TWISS}.}
Толстый \texttt{SBEND} описывается матрицами TRANSPORT первого и второго порядков~\cite{Brown1968,MADX}, куда входят и краевые матрицы диполя. Хроматичность \texttt{DQ} вычисляется по хроматическим функциям, т.\,е. по формулам вида \eqref{eq:xi-ham} (с флагом \texttt{CHROM} --- конечной разностью по $\delta$), и относится к $p_t$: $\xi=\beta_s\,\mathtt{DQ}$ \eqref{eq:xi-conv}. Член F присутствует даже при $E_1=E_2=0$, $\mathtt{FINT}=0$: краевая матрица второго порядка содержит кинематический член $\propto x'y$. Ненулевые \texttt{FINT}, \texttt{HGAP} добавляют поправку конечного края раздела~\ref{sec:edge} в приближении тонкого слоя.

\paragraph{PTC.}
В точной модели (\texttt{exact}) тело диполя интегрируется по полному корню \eqref{eq:H}, а жёсткий край всегда содержит кинематический краевой член, даже при выключенном \texttt{fringe}; \texttt{fringe=true} добавляет краевое поле. При \texttt{time=false} переменная --- $\delta_p$. Хроматичность и её высшие порядки получаются из нормальной формы с $\delta$ как параметром (\texttt{ptc\_normal}, \texttt{ptc\_twiss})~\cite{PTC}. Приближённая (expanded) модель теряет часть членов с $1/\rho$, а тонкие диполи (\texttt{MAKETHIN}) --- члены C, E и F; для хроматичности их использовать нельзя.

\subsection{OptiM}

OptiM~\cite{OptiM} вычисляет натуральную хроматичность по $\delta_p$ из линейной оптики секторного магнита с жёстким краем, включая краевой член F; результаты совпадают с MAD-X и PTC. Высших хроматичностей и краевого поля конечной длины в используемой модели нет.

\subsection{Python: точный гамильтониан}

Независимая проверка --- прямое интегрирование уравнений Гамильтона для \eqref{eq:H} в теле элементов, поиск замкнутой орбиты Ньютоном при фиксированном $\delta$ и вычисление \eqref{eq:xi-trace} по якобиану. Без краевых ударов результат воспроизводит модель <<только тело>> (BELL, COSY \texttt{FR 0}); с ударами \eqref{eq:kick} на торцах --- MAD-X, PTC и OptiM. Удары обратного знака не воспроизводят ни одну программу, что подтверждает знак \eqref{eq:kick}.
